\setcounterpageref{ex}{0}
\TaoPhieuKiemTra{01}

\setcounter{ex}{0}
\setcounter{bt}{0}

\textbf{\large I. PHẦN TRẮC NGHIỆM (3,0 điểm)}

{\itshape Khoanh tròn vào chữ cái đặt trước câu trả lời đúng. Mỗi câu đúng được 0,5 điểm.}

\begin{ex} % Câu 1
Số liền trước của số $60$ là
\choice
{$61$}
{\True $59$}
{$50$}
{$70$}
\loigiai{
Số liền trước của $60$ là: $60 - 1 = 59$.\\
\textbf{Đáp án đúng: B}
}
\end{ex}

\begin{ex} % Câu 2
Trong phép tính $35 + 24 = 59$, số $59$ được gọi là
\choice
{Số hạng}
{Số bị trừ}
{\True Tổng}
{Hiệu}
\loigiai{
Trong phép cộng, kết quả được gọi là Tổng.\\
\textbf{Đáp án đúng: C}
}
\end{ex}

\begin{ex} % Câu 3
Đơn vị ki-lô-gam được viết tắt là
\choice
{$\text{km}$}
{$\text{g}$}
{$\text{l}$}
{\True $\text{kg}$}
\loigiai{
Ki-lô-gam viết tắt là $\text{kg}$.\\
\textbf{Đáp án đúng: D}
}
\end{ex}

\begin{ex} % Câu 4
Đoạn thẳng $AB$ được giới hạn bởi mấy điểm?
\choice
{\True $2$ điểm ($A$ và $B$)}
{$1$ điểm}
{$3$ điểm}
{Không có điểm nào}
\loigiai{
Đoạn thẳng $AB$ bị giới hạn bởi 2 đầu mút $A$ và $B$.\\
\textbf{Đáp án đúng: A}
}
\end{ex}

\begin{ex} % Câu 5
Điền dấu thích hợp vào chỗ chấm: $9 + 6 \dots 8 + 8$.
\choice
{$>$}
{$=$}
{\True $<$}
{Không so sánh được}
\loigiai{
$9 + 6 = 15$; $8 + 8 = 16$. Vì $15 < 16$ nên điền dấu $<$.\\
\textbf{Đáp án đúng: C}
}
\end{ex}

\begin{ex} % Câu 6
Giá trị của biểu thức $9 + 5 - 6$ là
\choice
{$7$}
{\True $8$}
{$9$}
{$10$}
\loigiai{
$9 + 5 = 14$; $14 - 6 = 8$.\\
\textbf{Đáp án đúng: B}
}
\end{ex}

\vspace{4mm}
\textbf{\large II. PHẦN TỰ LUẬN (7,0 điểm)}

{\itshape Thí sinh trình bày chi tiết bài làm vào giấy kiểm tra.}

\begin{bt}[1,0 điểm]
Tính nhẩm:
\begin{multicols}{2}
\begin{enumerate}[a)]
  \item $9 + 5$
  \item $7 + 6$
  \item $13 - 6$
  \item $15 - 7$
\end{enumerate}
\end{multicols}
\loigiai{
\begin{itemize}
  \item $9 + 5 = 14$
  \item $7 + 6 = 13$
  \item $13 - 6 = 7$
  \item $15 - 7 = 8$
\end{itemize}
}
\end{bt}

\begin{bt}[1,5 điểm]
Đặt tính rồi tính:
\begin{multicols}{3}
\begin{enumerate}[a)]
  \item $25 + 35$
  \item $77 - 24$
  \item $43 + 54$
\end{enumerate}
\end{multicols}
\loigiai{
\begin{itemize}
  \item $25 + 35 = 60$
  \item $77 - 24 = 53$
  \item $43 + 54 = 97$
\end{itemize}
}
\end{bt}

\begin{bt}[1,0 điểm]
Tính giá trị biểu thức kèm đơn vị đo:
\begin{enumerate}[a)]
  \item $17\text{ kg} + 8\text{ kg} - 10\text{ kg}$
  \item $26\text{ }l - 12\text{ }l + 7\text{ }l$
\end{enumerate}
\loigiai{
\begin{enumerate}[a)]
  \item $17\text{ kg} + 8\text{ kg} - 10\text{ kg} = 25\text{ kg} - 10\text{ kg} = 15\text{ kg}$.
  \item $26\text{ }l - 12\text{ }l + 7\text{ }l = 14\text{ }l + 7\text{ }l = 21\text{ }l$.
\end{enumerate}
}
\end{bt}

\begin{bt}[1,5 điểm]
Bạn Lan có $16$ viên bi. Bạn Mai có nhiều hơn bạn Lan $5$ viên bi. Hỏi bạn Mai có bao nhiêu viên bi?
\loigiai{
Số viên bi bạn Mai có là:
\[ 16 + 5 = 21 \text{ (viên bi)} \]
\textbf{Đáp số:} $21$ viên bi.
}
\end{bt}

\begin{bt}[1,0 điểm]
Một thùng chứa đầy $36\text{ }l$ nước mắm. Buổi sáng cửa hàng đã bán ra $13\text{ }l$. Hỏi trong thùng còn lại bao nhiêu lít nước mắm?
\loigiai{
Trong thùng còn lại số lít nước mắm là:
\[ 36 - 13 = 23\text{ }l \]
\textbf{Đáp số:} $23\text{ }l$ nước mắm.
}
\end{bt}

\begin{bt}[1,0 điểm]
Tìm số có hai chữ số, biết chữ số hàng chục là số lớn nhất có một chữ số, chữ số hàng đơn vị kém chữ số hàng chục $4$ đơn vị.
\loigiai{
Chữ số hàng chục là $9$.\\
Chữ số hàng đơn vị là: $9 - 4 = 5$.\\
Vậy số cần tìm là $95$.\\
\textbf{Đáp số:} $95$.
}
\end{bt}
